\subsection{Adoption of Hybrid Mobile Architecture}
\begin{itemize}
    \item \textbf{Description:} The development strategy involves building a single web-based codebase wrapped in a native container (Capacitor), rather than developing two separate native applications (Kotlin for Android, Swift for iOS) or a pure Progressive Web App (PWA).
    \item \textbf{Rationale:} A pure PWA was discarded because it cannot reliably access high-frequency motion sensors in the background, which is critical for the "Road Quality Detection" requirement. Native development was discarded due to the high effort required for a small team. The Hybrid approach offers the necessary hardware access with a unified codebase, significantly reducing development and testing time.
\end{itemize}

\subsection{Usage of PostgreSQL with PostGIS}
\begin{itemize}
    \item \textbf{Description:} The system utilizes a Relational Database Management System (RDBMS) extended with the PostGIS spatial engine to store and query all geospatial data (Points, Linestrings, Polygons).
    \item \textbf{Rationale:} Generic databases are insufficient for the mathematical complexity of this project. PostGIS provides native, highly optimized functions for geospatial conflation and Nearest-Neighbor searches. Implementing these geometric algorithms from scratch in code would be error-prone and inefficient.
\end{itemize}

\subsection{Centralized Logic for Merging and Ranking}
\begin{itemize}
    \item \textbf{Description:} The core algorithms that determine the shape of public paths (Merging) and their value (Ranking) are executed exclusively on the Cloud Tier via Edge Functions, never on the client.
    \item \textbf{Rationale:} This decision ensures a \textbf{Single Source of Truth}. The merging algorithm requires a complete view of all existing paths in an area to make a decision, which a single client does not possess. Furthermore, calculating the Ranking on the server prevents malicious users from tampering with the client code to artificially inflate the score of their preferred paths.
\end{itemize}

\subsection{Direct-to-Database Communication Pattern}
\begin{itemize}
    \item \textbf{Description:} The architecture bypasses the traditional "API Gateway / Controller" layer for standard operations. Clients perform CRUD (Create, Read, Update, Delete) operations directly against the database's REST API.
    \item \textbf{Rationale:} The majority of user interactions (e.g., loading a profile, saving a track) are standard data operations that do not require complex business logic. Implementing a dedicated middle-layer for these simple tasks would introduce unnecessary latency and code maintenance ("Boilerplate"). Security is delegated to RLS (as described in 2.6.2), making the middle layer redundant for these tasks.
\end{itemize}